% Temario, objetivos y programa transcritos de las páginas 3--8 de
% source/byn.pdf. Las duraciones se conservan tal como aparecen en el original.
\chapter*{Temario, objetivos y programa}

\section*{1. Motivación \duration{1 h}}
Motivar al grupo respecto a la importancia de la protección radiológica y de la necesidad de que la seguridad constituya un hábito personal. Explicará la responsabilidad adquirida al manejar una fuente radioactiva.

\section*{2. Estructura atómica: terminología \duration{1 h}}
El alumno identificará definiciones de átomo, núcleo, protón, electrón, número atómico, número de masa y número de neutrones; aplicará correctamente esas definiciones y la simbología
\[
{}^{A}_{Z}X.
\]

\section*{3. Radioactividad y radiación \duration{1:30 h}}
El alumno podrá describir el fenómeno de la desintegración radioactiva y las características de las radiaciones alfa, beta y gamma. Explicará el concepto de vida media y podrá determinar la actividad actual de una fuente conociendo su vida media, actividad inicial y el tiempo transcurrido, empleando una gráfica o una tabla de decaimiento. Empleará correctamente las unidades de actividad: curie y becquerel.

\section*{4. Interacción de la radiación con la materia \duration{1 h}}
Definirá con sus palabras cada uno de los siguientes conceptos: ion, ionización, excitación, absorción, alcance o penetración, dispersión, radiación primaria, radiación dispersa y radiación transmitida.

\textbf{Práctica no. 1.} \emph{Absorción de las radiaciones beta y gamma.} \duration{1 h}

\section*{5. Exposición y dosis. Unidades \duration{1 h}}
El alumno explicará la diferencia entre exposición y dosis, exposición y rapidez de exposición, dosis y rapidez de dosis. Definirá roentgen, el rad y sus submúltiplos, así como las unidades de rapidez de exposición y de dosis. Describirá los conceptos de equivalente de dosis y rem, así como las unidades gray y sievert y sus equivalentes a rad y rem.

\section*{6. Efectos biológicos de la radiación \duration{2:30 h}}

\subsection*{6.1. Efectos somáticos y efectos genéticos. Radiosensibilidad}
Podrá señalar la diferencia entre los efectos somáticos y los efectos genéticos, así como la diferente radiosensibilidad de los distintos tipos de células, tejidos y órganos.

\subsection*{6.2. Relación dosis--efecto. Efecto de dosis aguda. Dosis crónica y efectos tardíos}
Describirá la diferencia entre dosis aguda y dosis crónica, distinguirá los valores de dosis que producen los distintos efectos, desde un eritema hasta una dosis letal, así como la diferencia entre una irradiación localizada y una irradiación a cuerpo entero.

\subsection*{6.3. Formas de exposición a la radiación. Radiación externa, contaminación y dosis interna}
El alumno distinguirá entre casos de exposición a radiación externa, ingestión, inhalación y dosis interna.

\section*{7. Medición de la radiación \duration{0:30 h}}
\textbf{Objetivo general.} El alumno manejará los diversos instrumentos de monitoreo, escogerá el adecuado para cada situación e interpretará correctamente las lecturas; quedará convencido de la ventaja de usar siempre un dosímetro personal, una alarma sonora y un instrumento de monitoreo, así como de cuidarlos adecuadamente.

\subsection*{7.1. Dispositivos de vigilancia personal}

\textbf{a) Dosímetros de bolsillo \duration{0:30 h}}

Entenderá el funcionamiento y manejará un dosímetro de bolsillo y su cargador, podrá verificar su funcionamiento y tomar las medidas pertinentes para su conservación. Interpretará correctamente las lecturas y conocerá sus ventajas y limitaciones.

\textbf{Práctica no. 2.} \emph{Manejo y uso del dosímetro de bolsillo de lectura directa} \duration{1:00 h}

\textbf{b) Dosímetros de película y termoluminiscentes \duration{0:30 h}}

El alumno reconocerá y podrá explicar la importancia del uso de un dosímetro que registre la exposición acumulada. Se percatará de la necesidad de adquirir el hábito de usarlo y cuidar de que su uso y conservación sean adecuados.

\textbf{c) Detector con alarma sonora \duration{0:30 h}}

El alumno usará un detector con alarma sonora, relacionará la señal audible con el índice de exposición correspondiente y reconocerá las ventajas de su uso continuo.

\textbf{Práctica no. 3.} \emph{Manejo y uso del detector con alarma sonora} \duration{1 h}

\subsection*{7.2. Instrumentos de monitoreo}
\textbf{Objetivo general.} El alumno manejará correctamente al menos uno de los siguientes instrumentos de monitoreo: con cámara de ionización, con detector Geiger o con detector de centelleo. Interpretará correctamente las lecturas tomando en consideración las limitaciones de cada uno. Podrá verificar la calibración del instrumento usando una fuente de referencia.

\textbf{a) Monitor con detector Geiger \duration{1 h}}

Lo manejará y explicará en qué casos es conveniente su uso. Describirá sus ventajas y desventajas frente a una cámara de ionización. Describirá el fenómeno de saturación y el uso del factor de calibración.

\textbf{Práctica no. 4.} \emph{Manejo y uso de un monitor con detector Geiger} \duration{1 h}

\textbf{Práctica no. 5.} \emph{Localización de fuentes con Geiger--Müller} \duration{1:30 h}

\section*{8. Filosofía de la protección radiológica \duration{2 h}}

\subsection*{8.1. Concepto de riesgo--beneficio. Concepto ALARA}
El alumno podrá percatarse de las diferencias entre el riesgo y el beneficio, en diversos tipos de actividades y, en especial, en la radiografía industrial. Explicará en forma simplificada el concepto ALARA.

\subsection*{8.2. Sistema de limitación de dosis}
Memorizará el valor del límite de equivalente de dosis para trabajadores ocupacionalmente expuestos y los valores de rapidez de exposición en mR/h y mR/semana que en promedio producirán dicho límite.

\section*{9. Principios básicos de protección contra radiación externa}

\subsection*{9.1. Distancia y tiempo como medios de protección \duration{2 h}}
El alumno calculará la rapidez de exposición a diferentes distancias de una fuente de actividad conocida. Calculará la exposición en términos de rapidez de exposición y el tiempo. Calculará el tiempo que puede permanecer una persona en un campo de radiación conocido, para recibir una exposición o una dosis predeterminada a distintas distancias de la fuente. Deberá quedar suficientemente familiarizado con este tipo de cálculos como para hacerlos sin dificultad en forma rutinaria.

\textbf{Práctica no. 6.} \emph{Aplicación de la ley del cuadrado inverso} \duration{1:30 h}

\subsection*{9.2. El blindaje como medio de protección \duration{1 h}}
El alumno podrá determinar la atenuación de un haz de radiación gamma al atravesar un blindaje de espesor y material conocidos empleando para ello una gráfica semilogarítmica y aplicando el concepto de una capa hemirreductora y capa decirreductora para determinar espesores de blindaje suficientes, a fin de producir la atenuación deseada.

\textbf{Práctica no. 7.} \emph{Atenuación de la radiación gamma, determinación de la capa hemirreductora utilizando diversos materiales (plomo, concreto, ladrillo, agua)} \duration{1 h}

\subsection*{9.3. Protección contra la contaminación externa y la ingestión o la inhalación \duration{0:30 h}}
El alumno reafirmará sus conceptos de contaminación y podrá explicar la diferencia entre contaminación externa, ingestión e inhalación, así como la forma en que pueden producirse y los medios para evitarla.

\subsection*{9.4. La protección radiológica durante el trabajo \duration{1 h}}
Identificará las situaciones de mayor riesgo durante el trabajo radiográfico. Describirá secuencialmente las medidas de protección que deben tomarse en una jornada de trabajo.

\section*{10. Equipos para radiografía industrial que utilizan material radiactivo}

\subsection*{10.1. Principios de la utilización de la fuente y descripción de los equipos asociados en el uso \duration{1 h}}

\subsection*{10.2. Procedimientos generales para el transporte y almacenamiento de material radiactivo \duration{1 h}}
El alumno reconocerá y podrá explicar la importancia de las medidas de seguridad física y radiológica que deben contemplarse en el transporte y durante el almacén del material radiactivo. Describirá las posibles consecuencias de pasar por alto dichas medidas.

\section*{11. Normas y reglamentos \duration{1:30 h}}
\textbf{Objetivo general.} Conocer las principales disposiciones reglamentarias aplicables y las funciones del organismo regulador en materia de seguridad radiológica.

\subsection*{11.1. Manual de seguridad radiológica}
El alumno dará una interpretación correcta (que le permita en el futuro aplicar adecuadamente) a cada una de las disposiciones del Manual Interno de Seguridad Radiológica. Distinguirá el contenido mínimo de un Manual de Seguridad Radiológica para Radiografía Industrial.

\subsection*{11.2. Estructura, organización y cadena de responsabilidad en materia de seguridad radiológica}
El alumno podrá canalizar adecuadamente sus necesidades de equipo de seguridad, información y ayuda en caso de incidente, dentro de la estructura jerárquica de la empresa.

\subsection*{11.3. Obligaciones del encargado de seguridad radiológica}
El alumno describirá las funciones más importantes del Encargado de Seguridad Radiológica.

\subsection*{11.4. La C.N.S.N.S., sus funciones y organización}
El alumno explicará cuáles son las funciones y las dependencias de la C.N.S.N.S. relacionadas con su área de trabajo, así como la forma de acceso a esas dependencias cuando requiera orientación en materia de seguridad radiológica o auxilio en caso de incidentes o accidentes radiológicos.

\section*{12. Análisis de incidentes y accidentes (casos reales) \duration{3 h}}
\textbf{Objetivo general.} Motivar mediante la exposición de incidentes y accidentes reales el análisis de las situaciones para buscar con calma y razonamiento lógico la solución que dé lugar a la exposición mínima, evitando así acciones precipitadas que pueden conducir a riesgos mayores.

\subsection*{12.1. Procedimiento de emergencia}
El alumno explicará cómo proceder en cada uno de los tipos de incidentes más comunes, en los que esté a su alcance la posibilidad de solución. Indicará a quién acudir cuando no disponga de los medios para resolver el caso, señalando las medidas o disposiciones inmediatas que debe tomar para evitar que se agrave la situación.

\subsection*{12.2. Informe de incidentes y accidentes}
El alumno indicará los tipos de incidentes que deberá informar a sus superiores y la forma en que lo hará. Deberá estar consciente de que todo accidente se informará de inmediato.

\subsection*{12.3. Plan de emergencia}

\textbf{Práctica no. 8.} \emph{Manejo de un equipo radiográfico, solución de fallas} \duration{2 h}

\textbf{Práctica no. 9.} \emph{Recuperación de fuentes, uso de manipuladores a distancia} \duration{2 h}

\textbf{Práctica no. 10.} \emph{Simulacro de un accidente} \duration{4 h}

\textbf{Evaluación final teórico-práctica.}

\textbf{Duración total del curso: 40 h (24 h teoría, 16 h práctica).}
